% ============================================================================
% 08_diagramas_colaboracion.tex — Cuerpo del entregable "Diagramas de
% Colaboración". Sin preámbulo: se incluye tanto desde
% 08_Diagramas_Colaboracion.tex (standalone) como, eventualmente, desde un
% documento maestro.
% ============================================================================
\section{Diagramas de Colaboración}

\subsection{Propósito y alcance}

Esta sección documenta, siguiendo la técnica descripta por Craig Larman en
\textit{UML y Patrones}, los \textbf{contratos de operación de sistema} y los
\textbf{diagramas de colaboración} (\textit{communication diagrams} en UML
2.5) correspondientes a las nueve operaciones de sistema identificadas en
\texttt{docs/IDS.md}: SO-01 a SO-09 (SO-06 a SO-08 cubren RF-03, RF-04 y
RF-06; SO-09 cubre RNF-03). Cada operación de sistema es un mensaje
que un actor le envía al sistema considerado como una caja negra durante el
diseño de la interacción; el diagrama de colaboración muestra cómo el
sistema, ya descompuesto en sus objetos de diseño reales (\textit{hooks} de
React, componentes, \textit{handlers} de FastAPI y el modelo de
reconocimiento), coopera para cumplir esa operación.

A diferencia de un diagrama de secuencia, un diagrama de colaboración no
tiene eje temporal ni líneas de vida: prioriza la topología de
colaboración entre objetos. Cada par de objetos que colabora está unido por
un único enlace estructural, dibujado como una línea recta \textbf{sin
punta de flecha}. Cada mensaje que viaja por ese enlace se dibuja como una
flechita corta y propia, paralela al enlace, apuntando siempre hacia el
objeto receptor, con su número jerárquico junto al texto del mensaje (1,
1.1, 1.2, 2, 2.1\ldots); cuando varios mensajes recorren el mismo enlace se
apilan sobre él, y los mensajes de sentido opuesto (por ejemplo, un mensaje
de vuelta genuino, como el evento asíncrono de un WebSocket) se dibujan del
lado contrario del enlace. Los valores de retorno de una llamada síncrona
simple no se dibujan como una flecha aparte: se integran en el propio
mensaje con la notación \texttt{resultado := operación(args)}, siguiendo la
convención de Larman de omitir flechas de retorno triviales. Un
automensaje (un objeto que se invoca a sí mismo) se dibuja como un pequeño
bucle sobre el propio objeto. Cada objeto se identifica como
\texttt{instancia~:~Clase} (o \texttt{:Clase} cuando no hace falta nombrar
la instancia), con el nombre subrayado, cuando corresponde a una instancia
real (por ejemplo, un componente de React o el modelo LSTM cargado en
memoria), o simplemente por su nombre cuando representa un módulo o una
colección de funciones del backend sin una clase propia (por ejemplo,
\texttt{capture\_validation} o \texttt{MEDICAL\_TEMPLATES}, el diccionario
de plantillas médicas para combinaciones frecuentes de señas), ya que el
backend de LSA Salud (\texttt{app.py}) no está organizado en clases salvo por el
modelo \texttt{SignLanguageLSTM}. Los nueve diagramas se dibujan en
draw.io (formato \texttt{.drawio} editable, generado desde
\texttt{docs/diagramas/colaboracion/gen\_drawio.py}) para poder expresar
esta notación de forma exacta; ver \texttt{docs/README.md} para el flujo de
edición y reexportación.

Para cada operación se presenta: el contrato de la operación (con sus
referencias cruzadas a caso de uso y requisito funcional), el diagrama de
colaboración, los patrones GRASP aplicados y una tabla que traza cada objeto
del diagrama a su archivo fuente real.

\subsection{SO-01: reconocerSeña(frame) --- DC-01}

\subsubsection{Contrato de operación}

\begin{table}[H]
\centering
\begin{tabularx}{\textwidth}{@{}p{4.4cm} X@{}}
\toprule
\textbf{Operación} & \texttt{reconocerSeña(frame)} \\
\textbf{Referencias cruzadas} & Realiza CU-01 (pasos de captura, reconocimiento
y armado de la secuencia); RF-01.1, RF-01.2, RF-01.3, RF-01.4. \\
\textbf{Precondiciones} & La cámara está autorizada y el \textit{hook} de
MediaPipe Tasks está emitiendo \textit{frames} (a 20 fps); el backend tiene el modelo
\texttt{SignLanguageLSTM} cargado; existe una conexión WebSocket activa hacia
\texttt{/ws/predict}. \\
\textbf{Postcondiciones} & El \textit{frame} quedó agregado al búfer
deslizante de la Estación del Paciente. Si el búfer alcanzó el múltiplo de
\texttt{PREDICT\_STRIDE\_FRAMES} configurado, se envió una predicción al
backend y, si la seña reconocida superó el umbral de confianza y no fue
``idle'', la palabra reconocida quedó agregada a la secuencia en curso
(\texttt{transcript}) de \texttt{useSentenceBuilder}; si en cambio fue
``idle'' con confianza $\geq$70\%, se invocó el cierre de la oración en
curso (SO-02). \\
\bottomrule
\end{tabularx}
\caption{Contrato de la operación SO-01 --- reconocerSeña(frame)}
\label{tab:col-contrato-so01}
\end{table}

\subsubsection{Diagrama de colaboración DC-01}

\diagrama{08_colaboracion_dc01}{\textwidth}{DC-01 --- Colaboración para reconocerSeña(frame)}{fig:col-dc01}

\subsubsection{Patrones GRASP aplicados}

\begin{itemize}[leftmargin=*, itemsep=2pt]
    \item \textbf{Controlador:} \texttt{handleFrame} en \texttt{PacientePage}
    actúa como controlador de la operación de sistema: recibe el evento de
    captura desde el \textit{hook} de MediaPipe y coordina al búfer, al
    constructor de oraciones y al \textit{hook} de WebSocket, sin contener
    él mismo la lógica de reconocimiento.
    \item \textbf{Experto en información:} el modelo \texttt{SignLanguageLSTM}
    es quien tiene el conocimiento (los pesos entrenados) para clasificar el
    tensor de entrada, por lo que es el responsable de producir la
    predicción; ningún otro objeto duplica esa responsabilidad.
    \item \textbf{Bajo acoplamiento:} \texttt{useSignWebSocket} aísla a
    \texttt{PacientePage} del protocolo de transporte (JSON o binario) y de
    los detalles de reconexión del WebSocket; \texttt{PacientePage} solo
    conoce \texttt{sendFrames} y un \textit{callback} \texttt{onMessage}.
    \item \textbf{Indirección:} el \textit{handler}
    \texttt{websocket\_predict} media entre el frontend y el modelo
    \texttt{SignLanguageLSTM}: el frontend nunca invoca al modelo
    directamente, sino a través de este intermediario, lo que permite
    cambiar el modelo o el preprocesamiento sin afectar al cliente.
\end{itemize}

\subsubsection{Objetos del diagrama y su ubicación en el código}

\begin{table}[H]
\centering
\begin{tabularx}{\textwidth}{@{}l X@{}}
\toprule
\textbf{Objeto} & \textbf{Ubicación real} \\
\midrule
\texttt{useMediaPipeTasks} & Hook $\rightarrow$ \texttt{web/hooks/useMediaPipeTasks.ts}; extrae landmarks de manos, pose y rostro (\textit{blendshapes}) en el navegador, a 20 fps fijos, y expone el contrato \texttt{onFrame} \\
\texttt{handleFrame : PacientePage} & Componente $\rightarrow$ \texttt{web/app/paciente/page.tsx} \\
\texttt{recordedFrames} & Ref (búfer) $\rightarrow$ \texttt{web/app/paciente/page.tsx} \\
\texttt{useSignWebSocket} & Hook $\rightarrow$ \texttt{web/hooks/useSignWebSocket.ts} \\
\texttt{websocket\_predict} & Handler FastAPI $\rightarrow$ \texttt{app.py} \\
\texttt{model : SignLanguageLSTM} & Instancia del modelo $\rightarrow$ \texttt{src/model/lstm.py} \\
\texttt{TopPredictions} & Componente $\rightarrow$ \texttt{web/components/*/TopPredictions.tsx} \\
\texttt{useSentenceBuilder} & Hook $\rightarrow$ \texttt{web/hooks/useSentenceBuilder.ts} \\
\bottomrule
\end{tabularx}
\caption{Trazabilidad de objetos --- DC-01}
\label{tab:col-objetos-dc01}
\end{table}

\subsection{SO-02: cerrarOracionYTraducir(secuencia) --- DC-02}

\subsubsection{Contrato de operación}

\begin{table}[H]
\centering
\begin{tabularx}{\textwidth}{@{}p{4.4cm} X@{}}
\toprule
\textbf{Operación} & \texttt{cerrarOracionYTraducir(secuencia)} \\
\textbf{Referencias cruzadas} & Realiza CU-01 (cierre de secuencia) y CU-02
completo; RF-01.4, RF-01.5. \\
\textbf{Precondiciones} & La secuencia de palabras reconocidas
(\texttt{transcript}) no está vacía; el modelo predijo la clase de reposo
(\texttt{idle}) con una confianza $\geq$70\% y no hay otra oración en
procesamiento. \\
\textbf{Postcondiciones} & La secuencia en curso quedó cerrada y vaciada
(\texttt{transcriptRef} vuelve a estar vacío). Se obtuvo una oración natural
en español --- por plantilla médica, por caché de traducciones o por el
modelo de lenguaje local --- y, si se obtuvo una oración válida, el mensaje
del paciente quedó agregado a la conversación compartida a través de
\texttt{useChatChannel}. \\
\bottomrule
\end{tabularx}
\caption{Contrato de la operación SO-02 --- cerrarOracionYTraducir(secuencia)}
\label{tab:col-contrato-so02}
\end{table}

\subsubsection{Diagrama de colaboración DC-02}

\diagrama{08_colaboracion_dc02}{0.8\textwidth}{DC-02 --- Colaboración para cerrarOracionYTraducir(secuencia)}{fig:col-dc02}

\subsubsection{Patrones GRASP aplicados}

\begin{itemize}[leftmargin=*, itemsep=2pt]
    \item \textbf{Experto en información:} \texttt{useSentenceBuilder} decide
    cuándo cerrar la oración porque es quien tiene el \textit{transcript}
    en curso y el indicador de oración en procesamiento, y quien recibe la
    predicción de reposo que dispara el cierre; ningún otro objeto podría
    tomar esa decisión sin duplicar ese estado.
    \item \textbf{Controlador:} \texttt{process\_sentence} en \texttt{app.py}
    coordina la resolución de la traducción (plantilla, caché, modelo de
    lenguaje) como punto único de entrada del backend para esta operación.
    \item \textbf{Indirección:} \texttt{MEDICAL\_TEMPLATES} y
    \texttt{\_translation\_cache} median entre la secuencia de palabras y el
    modelo de lenguaje externo, evitando invocar a Ollama cuando ya existe
    una traducción conocida o reciente.
    \item \textbf{Bajo acoplamiento:} \texttt{useSentenceBuilder} no conoce
    cómo se traduce la secuencia (plantilla, caché o LLM); solo conoce el
    contrato HTTP de \texttt{/process\_sentence}, lo que permite cambiar la
    estrategia de traducción en el backend sin tocar el frontend.
\end{itemize}

\subsubsection{Objetos del diagrama y su ubicación en el código}

\begin{table}[H]
\centering
\begin{tabularx}{\textwidth}{@{}l X@{}}
\toprule
\textbf{Objeto} & \textbf{Ubicación real} \\
\midrule
\texttt{useSentenceBuilder} & Hook $\rightarrow$ \texttt{web/hooks/useSentenceBuilder.ts} \\
\texttt{process\_sentence} & Handler FastAPI $\rightarrow$ \texttt{app.py} \\
\texttt{MEDICAL\_TEMPLATES} & Diccionario de módulo $\rightarrow$ \texttt{app.py} \\
\texttt{\_translation\_cache} & Caché LRU en memoria $\rightarrow$ \texttt{app.py} \\
\texttt{Ollama (qwen2:1.5b)} & Proceso local externo $\rightarrow$ API HTTP en \texttt{localhost:11434} \\
\texttt{useChatChannel (Paciente)} & Hook $\rightarrow$ \texttt{web/hooks/useChatChannel.ts} \\
\bottomrule
\end{tabularx}
\caption{Trazabilidad de objetos --- DC-02}
\label{tab:col-objetos-dc02}
\end{table}

\subsection{SO-03: sintetizarVoz(oracion) --- DC-03}

\subsubsection{Contrato de operación}

\begin{table}[H]
\centering
\begin{tabularx}{\textwidth}{@{}p{4.4cm} X@{}}
\toprule
\textbf{Operación} & \texttt{sintetizarVoz(oracion)} \\
\textbf{Referencias cruzadas} & Realiza CU-03 completo; RF-01.6. \\
\textbf{Precondiciones} & Existe una oración traducida no vacía
(resultado de SO-02); el motor Piper TTS (\texttt{piper\_voice}) puede estar
cargado o no. \\
\textbf{Postcondiciones} & Si el motor de síntesis está disponible, el audio
de la oración quedó disponible en formato \textit{base64} --- reutilizado
desde \texttt{\_audio\_cache} si ya existía, o generado y agregado a la
caché en caso contrario --- y se programó su reproducción en el navegador
tras la ventana de exhibición previa. Si el motor no está disponible, el
mensaje queda entregado sin audio, sin interrumpir el flujo de la
conversación. \\
\bottomrule
\end{tabularx}
\caption{Contrato de la operación SO-03 --- sintetizarVoz(oracion)}
\label{tab:col-contrato-so03}
\end{table}

\subsubsection{Diagrama de colaboración DC-03}

\diagrama{08_colaboracion_dc03}{\textwidth}{DC-03 --- Colaboración para sintetizarVoz(oracion)}{fig:col-dc03}

\subsubsection{Patrones GRASP aplicados}

\begin{itemize}[leftmargin=*, itemsep=2pt]
    \item \textbf{Experto en información:} \texttt{piper\_voice} es el único
    objeto con el conocimiento (modelo de voz \texttt{es\_AR-daniela-high})
    necesario para sintetizar el audio, por lo que es el responsable de
    generarlo.
    \item \textbf{Indirección:} \texttt{\_audio\_cache} media entre
    \texttt{process\_sentence} y \texttt{piper\_voice}, evitando resintetizar
    audio ya generado para la misma oración.
    \item \textbf{Bajo acoplamiento:} \texttt{useSentenceBuilder} no conoce
    el motor de síntesis ni el formato interno del audio; solo recibe
    \texttt{audio\_base64} y decide cuándo reproducirlo, tras la ventana de
    exhibición configurada por \texttt{DISPLAY\_WINDOW\_MS}.
    \item \textbf{Alta cohesión:} toda la responsabilidad de obtener el
    audio (consultar caché, sintetizar, cachear) queda concentrada en
    \texttt{process\_sentence}, sin dispersarse hacia el frontend.
\end{itemize}

\subsubsection{Objetos del diagrama y su ubicación en el código}

\begin{table}[H]
\centering
\begin{tabularx}{\textwidth}{@{}l X@{}}
\toprule
\textbf{Objeto} & \textbf{Ubicación real} \\
\midrule
\texttt{process\_sentence} & Handler FastAPI $\rightarrow$ \texttt{app.py} \\
\texttt{\_audio\_cache} & Caché LRU en memoria (hasta 50 oraciones) $\rightarrow$ \texttt{app.py} \\
\texttt{piper\_voice : PiperVoice} & Instancia del motor TTS $\rightarrow$ \texttt{app.py} \\
\texttt{useSentenceBuilder} & Hook $\rightarrow$ \texttt{web/hooks/useSentenceBuilder.ts} \\
\texttt{Audio} & API nativa del navegador, invocada desde \texttt{web/hooks/useSentenceBuilder.ts} \\
\bottomrule
\end{tabularx}
\caption{Trazabilidad de objetos --- DC-03}
\label{tab:col-objetos-dc03}
\end{table}

\subsection{SO-04: enviarRespuesta(mensaje) --- DC-04}

\subsubsection{Contrato de operación}

\begin{table}[H]
\centering
\begin{tabularx}{\textwidth}{@{}p{4.4cm} X@{}}
\toprule
\textbf{Operación} & \texttt{enviarRespuesta(mensaje)} \\
\textbf{Referencias cruzadas} & Realiza CU-04 (sincronización por
BroadcastChannel); RF-01.7, RF-02.1. \\
\textbf{Precondiciones} & La Estación de Recepción está abierta; el
personal de recepción o enfermería redactó un texto no vacío. \\
\textbf{Postcondiciones} & El mensaje del personal quedó agregado a
\texttt{chatHistory} en ambas estaciones (Recepción y Paciente), de forma
sincronizada a través del canal \texttt{nexora-chat}, y visible para el
paciente en su propia pantalla. \\
\bottomrule
\end{tabularx}
\caption{Contrato de la operación SO-04 --- enviarRespuesta(mensaje)}
\label{tab:col-contrato-so04}
\end{table}

\subsubsection{Diagrama de colaboración DC-04}

\diagrama{08_colaboracion_dc04}{0.85\textwidth}{DC-04 --- Colaboración para enviarRespuesta(mensaje)}{fig:col-dc04}

\subsubsection{Patrones GRASP aplicados}

\begin{itemize}[leftmargin=*, itemsep=2pt]
    \item \textbf{Controlador:} \texttt{RecepcionPage} recibe el evento de
    envío desde \texttt{ReceptionControls} y lo delega a
    \texttt{useChatChannel}, sin conocer los detalles de sincronización
    entre estaciones.
    \item \textbf{Indirección:} la \texttt{BroadcastChannel} nativa del
    navegador media entre las dos instancias de \texttt{useChatChannel}
    (Recepción y Paciente), que nunca se referencian directamente entre sí.
    \item \textbf{Bajo acoplamiento:} \texttt{ReceptionControls} solo conoce
    un \textit{callback} \texttt{onSend(texto)}; no sabe nada sobre el canal
    de sincronización ni sobre la estación del paciente.
    \item \textbf{Alta cohesión:} \texttt{useChatChannel} concentra toda la
    responsabilidad de publicar y recibir mensajes por el canal, tanto del
    lado de Recepción como del lado del Paciente (misma implementación,
    instanciada en cada estación).
\end{itemize}

\subsubsection{Objetos del diagrama y su ubicación en el código}

\begin{table}[H]
\centering
\begin{tabularx}{\textwidth}{@{}l X@{}}
\toprule
\textbf{Objeto} & \textbf{Ubicación real} \\
\midrule
\texttt{ReceptionControls} & Componente $\rightarrow$ \texttt{web/components/reception/ReceptionControls.tsx} \\
\texttt{RecepcionPage} & Componente $\rightarrow$ \texttt{web/app/recepcion/page.tsx} \\
\texttt{useChatChannel (Recepción)} & Hook $\rightarrow$ \texttt{web/hooks/useChatChannel.ts} \\
\texttt{BroadcastChannel} & API nativa del navegador (canal \texttt{nexora-chat}) \\
\texttt{useChatChannel (Paciente)} & Hook $\rightarrow$ \texttt{web/hooks/useChatChannel.ts} \\
\texttt{ChatBubble (Paciente)} & Componente $\rightarrow$ \texttt{web/components/chat/ChatBubble.tsx} \\
\bottomrule
\end{tabularx}
\caption{Trazabilidad de objetos --- DC-04}
\label{tab:col-objetos-dc04}
\end{table}

\subsection{SO-05: capturarMuestra(frames, etiqueta) --- DC-05}

\subsubsection{Contrato de operación}

\begin{table}[H]
\centering
\begin{tabularx}{\textwidth}{@{}p{4.4cm} X@{}}
\toprule
\textbf{Operación} & \texttt{capturarMuestra(frames, etiqueta)} \\
\textbf{Referencias cruzadas} & Realiza CU-06 (grabación de una repetición);
RF-T-03, RF-T-04. \\
\textbf{Precondiciones} & La Estación del Paciente está en Modo Captura; se
seleccionó o creó una etiqueta de seña; la cuenta regresiva de la repetición
ya finalizó. \\
\textbf{Postcondiciones} & Si los datos son válidos (etiqueta bien formada,
cantidad de \textit{frames} y de \textit{blendshapes} consistente), la
muestra quedó persistida siempre como dos archivos \texttt{.npy}, con los
landmarks en \texttt{dataset/custom/} y los \textit{blendshapes} en
\texttt{dataset/custom\_blendshapes/}, ambos asociados a la etiqueta y con la
misma marca de tiempo, y el conteo de muestras por etiqueta quedó actualizado. Si los
datos no son válidos, no se guardó ningún archivo y el error quedó
informado al Modo Captura. \\
\bottomrule
\end{tabularx}
\caption{Contrato de la operación SO-05 --- capturarMuestra(frames, etiqueta)}
\label{tab:col-contrato-so05}
\end{table}

\subsubsection{Diagrama de colaboración DC-05}

\diagrama{08_colaboracion_dc05}{\textwidth}{DC-05 --- Colaboración para capturarMuestra(frames, etiqueta)}{fig:col-dc05}

\subsubsection{Patrones GRASP aplicados}

\begin{itemize}[leftmargin=*, itemsep=2pt]
    \item \textbf{Controlador:} \texttt{CaptureMode} coordina la secuencia de
    grabación (cuenta regresiva, acumulación de \textit{frames}, envío) como
    punto único de entrada de la operación en el frontend.
    \item \textbf{Experto en información:} \texttt{capture\_validation}
    concentra el conocimiento de las reglas de validación (formato de
    etiqueta, consistencia de la cantidad de \textit{frames} y de
    \textit{blendshapes}) y es quien debe decidir si los datos son
    aceptables, no el \textit{handler} \texttt{capture}.
    \item \textbf{Creador:} el \textit{handler} \texttt{capture} es
    responsable de crear el archivo \texttt{.npy} persistido, porque es
    quien contiene los datos ya validados y quien conoce la ruta de
    destino.
    \item \textbf{Bajo acoplamiento:} \texttt{CaptureMode} no conoce las
    reglas de validación ni el formato de persistencia en disco; solo conoce
    el contrato HTTP de \texttt{/capture} y su respuesta.
\end{itemize}

\subsubsection{Objetos del diagrama y su ubicación en el código}

\begin{table}[H]
\centering
\begin{tabularx}{\textwidth}{@{}l X@{}}
\toprule
\textbf{Objeto} & \textbf{Ubicación real} \\
\midrule
\texttt{CaptureMode} & Componente $\rightarrow$ \texttt{web/components/capture/CaptureMode.tsx} \\
\texttt{captureHandlerRef} & Ref (manejador de \textit{frames} en modo captura) $\rightarrow$ \texttt{web/components/capture/CaptureMode.tsx} \\
\texttt{capture} & Handler FastAPI $\rightarrow$ \texttt{app.py} \\
\texttt{capture\_validation} & Módulo de funciones de validación $\rightarrow$ \texttt{capture\_validation.py} \\
\texttt{dataset/custom/<etiqueta>/*.npy} & Archivo de landmarks persistido en disco $\rightarrow$ \texttt{dataset/custom/} \\
\texttt{dataset/custom\_blendshapes/<etiqueta>/*.npy} & Archivo de \textit{blendshapes} persistido en disco $\rightarrow$ \texttt{dataset/custom\_blendshapes/} \\
\bottomrule
\end{tabularx}
\caption{Trazabilidad de objetos --- DC-05}
\label{tab:col-objetos-dc05}
\end{table}

\subsection{SO-06: otorgarConsentimiento(paciente) --- DC-06}
\textit{Cubre el requisito RF-03.}

\subsubsection{Contrato de operación}

\begin{table}[H]
\centering
\begin{tabularx}{\textwidth}{@{}p{4.4cm} X@{}}
\toprule
\textbf{Operación} & \texttt{otorgarConsentimiento(paciente)} \\
\textbf{Referencias cruzadas} & Realiza CU-09 completo; RF-03.1. \\
\textbf{Precondiciones} & \texttt{PacientePage} muestra el diálogo de
consentimiento (P-05); la cámara todavía no fue activada
(\texttt{consentState !== "accepted"}). \\
\textbf{Postcondiciones} & El consentimiento quedó registrado como evento
de auditoría anónimo (\texttt{type: "consent"}) en la tabla
\texttt{audit\_event} de la base de datos SQLite
(\texttt{data/lsa\_salud.db}), y --- si fue aceptado --- \texttt{PacientePage}
pasó \texttt{enabled: true} al hook de landmarks (\texttt{useMediaPipeTasks}),
habilitando la cámara y la Pantalla P-02. \\
\bottomrule
\end{tabularx}
\caption{Contrato de la operación SO-06 --- otorgarConsentimiento(paciente)}
\label{tab:col-contrato-so06}
\end{table}

\subsubsection{Diagrama de colaboración DC-06}

\diagrama{08_colaboracion_dc06}{\textwidth}{DC-06 --- Colaboración para otorgarConsentimiento(paciente)}{fig:col-dc06}

\subsubsection{Patrones GRASP aplicados}

\begin{itemize}[leftmargin=*, itemsep=2pt]
    \item \textbf{Controlador:} \texttt{PacientePage} recibe el clic del
    diálogo de consentimiento y coordina tanto el envío del evento de
    auditoría como la habilitación del hook de captura, como único punto
    de entrada de la operación en el frontend.
    \item \textbf{Experto en información:} \texttt{\_log\_audit} (dentro de
    \texttt{app.py}) delega en \texttt{storage.log\_audit\_event} el
    conocimiento del formato de persistencia de los eventos anónimos, en
    vez de que cada \textit{handler} sepa escribir su propio registro.
    \item \textbf{Indirección:} \texttt{sendAudit} (\texttt{audit.ts}) media
    entre \texttt{PacientePage} y el endpoint \texttt{POST /audit}, igual
    que \texttt{capture\_validation} media entre \texttt{capture} y el
    sistema de archivos en DC-05.
    \item \textbf{Bajo acoplamiento:} el diálogo de consentimiento no conoce
    el contrato HTTP ni el formato de persistencia; \texttt{PacientePage}
    solo invoca \texttt{sendAudit(\ldots)}, que es \textit{fire-and-forget}
    (no bloquea la habilitación de la cámara).
\end{itemize}

\subsubsection{Objetos del diagrama y su ubicación en el código}

\begin{table}[H]
\centering
\begin{tabularx}{\textwidth}{@{}l X@{}}
\toprule
\textbf{Objeto} & \textbf{Ubicación real} \\
\midrule
\texttt{PacientePage (diálogo de consentimiento)} & Componente $\rightarrow$ \texttt{web/app/paciente/page.tsx} \\
\texttt{sendAudit()} & Función $\rightarrow$ \texttt{web/lib/audit.ts} \\
\texttt{audit\_event} & Handler FastAPI (\texttt{POST /audit}) $\rightarrow$ \texttt{app.py} \\
\texttt{:useMediaPipeTasks} & Hook $\rightarrow$ \texttt{web/hooks/useMediaPipeTasks.ts} \\
\texttt{storage.log\_audit\_event()} & Función de persistencia relacional $\rightarrow$ \texttt{storage.py} \\
\texttt{data/lsa\_salud.db} & Base de datos SQLite persistida en disco $\rightarrow$ \texttt{data/} \\
\bottomrule
\end{tabularx}
\caption{Trazabilidad de objetos --- DC-06}
\label{tab:col-objetos-dc06}
\end{table}

\subsection{SO-07: registrarEventoAuditoria(operacion) --- DC-07}
\textit{Cubre el requisito RF-04.}

\subsubsection{Contrato de operación}

\begin{table}[H]
\centering
\begin{tabularx}{\textwidth}{@{}p{4.4cm} X@{}}
\toprule
\textbf{Operación} & \texttt{registrarEventoAuditoria(operacion)} \\
\textbf{Referencias cruzadas} & Realiza CU-11, incluido por CU-01 y CU-04; RF-04.1. \\
\textbf{Precondiciones} & Un \textit{handler} del backend (\texttt{predict},
\texttt{websocket\_predict}, \texttt{process\_sentence},
\texttt{report\_translation} o \texttt{audit\_event}) completó una
operación de sistema, o recibió un evento anónimo del frontend. \\
\textbf{Postcondiciones} & Un evento anónimo (tipo, marca temporal UTC y
campos propios del tipo) quedó registrado en la tabla \texttt{audit\_event}
de la base de datos SQLite (\texttt{data/lsa\_salud.db}), sin datos
identificatorios ni imágenes, video o texto de la conversación. \\
\bottomrule
\end{tabularx}
\caption{Contrato de la operación SO-07 --- registrarEventoAuditoria(operacion)}
\label{tab:col-contrato-so07}
\end{table}

\subsubsection{Diagrama de colaboración DC-07}

\diagrama{08_colaboracion_dc07}{0.9\textwidth}{DC-07 --- Colaboración para registrarEventoAuditoria(operacion)}{fig:col-dc07}

\subsubsection{Patrones GRASP aplicados}

\begin{itemize}[leftmargin=*, itemsep=2pt]
    \item \textbf{Experto en información:} \texttt{\_log\_audit} y el
    módulo de persistencia \texttt{storage.py} concentran el conocimiento
    del esquema de eventos de auditoría (tipo, marca temporal y métricas),
    asegurando la consistencia transaccional.
    \item \textbf{Bajo acoplamiento:} ningún \textit{handler} necesita
    conocer las sentencias SQL ni la ruta física de la base de datos: cada
    uno invoca \texttt{\_log\_audit(tipo, **campos)} con sus propios campos.
    \item \textbf{Fabricación pura:} \texttt{storage.log\_audit\_event}
    concentra la responsabilidad de inserción estructurada en SQLite
    (\texttt{data/lsa\_salud.db}), aislando la lógica de negocio de FastAPI
    del acceso a datos.
\end{itemize}

\subsubsection{Objetos del diagrama y su ubicación en el código}

\begin{table}[H]
\centering
\begin{tabularx}{\textwidth}{@{}l X@{}}
\toprule
\textbf{Objeto} & \textbf{Ubicación real} \\
\midrule
\texttt{handler} & Cualquiera de los \textit{handlers} FastAPI que completan una operación de sistema (\texttt{predict}, \texttt{websocket\_predict}, \texttt{process\_sentence}, \texttt{report}, \texttt{audit\_event}) $\rightarrow$ \texttt{app.py} \\
\texttt{\_log\_audit()} & Función envoltura $\rightarrow$ \texttt{app.py} \\
\texttt{storage.log\_audit\_event()} & Función de persistencia relacional $\rightarrow$ \texttt{storage.py} \\
\texttt{data/lsa\_salud.db} & Base de datos SQLite persistida en disco $\rightarrow$ \texttt{data/} \\
\bottomrule
\end{tabularx}
\caption{Trazabilidad de objetos --- DC-07}
\label{tab:col-objetos-dc07}
\end{table}

\subsection{SO-08: reportarTraduccion(mensaje) --- DC-08}
\textit{Cubre el requisito RF-06.}

\subsubsection{Contrato de operación}

\begin{table}[H]
\centering
\begin{tabularx}{\textwidth}{@{}p{4.4cm} X@{}}
\toprule
\textbf{Operación} & \texttt{reportarTraduccion(mensaje)} \\
\textbf{Referencias cruzadas} & Realiza CU-10, extiende a CU-01 y CU-04; RF-06.1. \\
\textbf{Precondiciones} & Existe un mensaje del paciente visible en el
panel de conversación de P-02 o P-03; el actor hizo clic en su botón
``Reportar traducción''. \\
\textbf{Postcondiciones} & El mensaje y su reporte (texto, estación,
palabras reconocidas y marca temporal) quedaron persistidos en las tablas
\texttt{reported\_message}, \texttt{recognized\_word} y
\texttt{translation\_report} de SQLite (\texttt{data/lsa\_salud.db}), sin
datos identificatorios del paciente; se registró un evento de auditoría
\texttt{report} en \texttt{audit\_event} (ver DC-07); el componente
\texttt{ChatBubble} emitió el evento \texttt{REPORTED} mediante
\texttt{BroadcastChannel} sincronizando el estado a la otra cara de la
estación; el texto del mensaje quedó marcado como tachado con advertencia
operativa y botón en estado permanente ``Traducción reportada'' en ambas
pantallas. \\
\bottomrule
\end{tabularx}
\caption{Contrato de la operación SO-08 --- reportarTraduccion(mensaje)}
\label{tab:col-contrato-so08}
\end{table}

\subsubsection{Diagrama de colaboración DC-08}

\diagrama{08_colaboracion_dc08}{0.9\textwidth}{DC-08 --- Colaboración para reportarTraduccion(mensaje)}{fig:col-dc08}

\subsubsection{Patrones GRASP aplicados}

\begin{itemize}[leftmargin=*, itemsep=2pt]
    \item \textbf{Bajo acoplamiento:} \texttt{ChatBubble} llama a
    \texttt{POST /report} con el payload estructurado y, ante la respuesta
    exitosa, notifica a su hook \texttt{useChatChannel} para propagar el
    evento por \texttt{BroadcastChannel}. No conoce detalles del motor de
    base de datos ni del esquema SQL.
    \item \textbf{Experto en información:} \texttt{validate\_report\_payload}
    valida las restricciones de longitud, estación y palabras reconocidas.
    El módulo \texttt{storage.py} garantiza la unicidad del mensaje
    reportado (\texttt{UNIQUE (client\_message\_id, text)}) y asocia
    múltiples reportes si ambas estaciones lo señalan.
    \item \textbf{Fabricación pura:} \texttt{storage.save\_report}
    encapsula la transacción relacional ACID (inserción de
    \texttt{reported\_message}, \texttt{recognized\_word} y
    \texttt{translation\_report}), protegiendo la base de datos contra
    inconsistencias.
    \item \textbf{Controlador:} el handler \texttt{report\_translation} en
    \texttt{app.py} coordina la validación, la persistencia en base de
    datos y el registro del evento de auditoría.
\end{itemize}

\subsubsection{Objetos del diagrama y su ubicación en el código}

\begin{table}[H]
\centering
\begin{tabularx}{\textwidth}{@{}l X@{}}
\toprule
\textbf{Objeto} & \textbf{Ubicación real} \\
\midrule
\texttt{ChatBubble} & Componente de interfaz $\rightarrow$ \texttt{web/components/chat/ChatBubble.tsx} \\
\texttt{useChatChannel} & Hook con \texttt{BroadcastChannel("nexora-chat")} $\rightarrow$ \texttt{web/hooks/useChatChannel.ts} \\
\texttt{report\_translation} & Handler FastAPI (\texttt{POST /report}) $\rightarrow$ \texttt{app.py} \\
\texttt{validate\_report\_payload()} & Función de validación $\rightarrow$ \texttt{capture\_validation.py} \\
\texttt{storage.save\_report()} & Servicio de persistencia SQLite $\rightarrow$ \texttt{storage.py} \\
\texttt{data/lsa\_salud.db} & Base de datos relacional SQLite $\rightarrow$ \texttt{data/} \\
\bottomrule
\end{tabularx}
\caption{Trazabilidad de objetos --- DC-08}
\label{tab:col-objetos-dc08}
\end{table}

\subsection{SO-09: limpiarConversacion() --- DC-09}
\textit{Cubre el requisito RNF-03.}

\subsubsection{Contrato de operación}

\begin{table}[H]
\centering
\begin{tabularx}{\textwidth}{@{}p{4.4cm} X@{}}
\toprule
\textbf{Operación} & \texttt{limpiarConversacion()} \\
\textbf{Referencias cruzadas} & Realiza CU-12 completo; RNF-03;
RF-04.1. \\
\textbf{Precondiciones} & Hay una atención en curso o recién finalizada en
la estación (la conversación no está vacía en al menos una de sus dos
caras). \\
\textbf{Postcondiciones} & La conversación quedó vacía en las dos caras de
la estación; se registró un evento de auditoría anónimo
(\texttt{type: "fin\_atencion"}) en la tabla \texttt{audit\_event} de la
base de datos SQLite (\texttt{data/lsa\_salud.db}); no queda ningún dato de
la sesión atendida. \\
\bottomrule
\end{tabularx}
\caption{Contrato de la operación SO-09 --- limpiarConversacion()}
\label{tab:col-contrato-so09}
\end{table}

\subsubsection{Diagrama de colaboración DC-09}

\diagrama{08_colaboracion_dc09}{\textwidth}{DC-09 --- Colaboración para limpiarConversacion()}{fig:col-dc09}

\subsubsection{Patrones GRASP aplicados}

\begin{itemize}[leftmargin=*, itemsep=2pt]
    \item \textbf{Controlador:} \texttt{RecepcionPage} centraliza el manejo
    del evento del botón ``Finalizar atención'' y lo delega en
    \texttt{useChatChannel}, sin conocer los detalles de sincronización
    entre las dos caras de la estación.
    \item \textbf{Bajo acoplamiento:} \texttt{useChatChannel} encapsula el
    mecanismo de sincronización entre las dos caras; ni \texttt{RecepcionPage}
    ni la otra cara de la estación conocen los detalles del canal.
    \item \textbf{Variación protegida:} el mecanismo de sincronización
    (canal de difusión entre las dos caras) queda aislado detrás de
    \texttt{useChatChannel}, igual que en DC-04.
    \item \textbf{Fabricación pura:} \texttt{audit.ts} es un módulo dedicado
    exclusivamente a enviar eventos de auditoría, reutilizado sin cambios
    desde DC-06 y DC-07.
\end{itemize}

\subsubsection{Objetos del diagrama y su ubicación en el código}

\begin{table}[H]
\centering
\begin{tabularx}{\textwidth}{@{}l X@{}}
\toprule
\textbf{Objeto} & \textbf{Ubicación real} \\
\midrule
\texttt{RecepcionPage} & Componente con el botón ``Finalizar atención'' $\rightarrow$ \texttt{web/app/recepcion/page.tsx} \\
\texttt{useChatChannel (Recepción)} & Hook (\texttt{clearHistory()}) $\rightarrow$ \texttt{web/hooks/useChatChannel.ts} \\
\texttt{BroadcastChannel} & API nativa del navegador (canal de sincronización entre las dos caras de la estación) \\
\texttt{useChatChannel (Paciente)} & Hook $\rightarrow$ \texttt{web/hooks/useChatChannel.ts} \\
\texttt{sendAudit()} & Función $\rightarrow$ \texttt{web/lib/audit.ts} \\
\texttt{audit\_event} & Handler FastAPI (\texttt{POST /audit}) $\rightarrow$ \texttt{app.py} \\
\texttt{storage.log\_audit\_event()} & Función de persistencia relacional $\rightarrow$ \texttt{storage.py} \\
\texttt{data/lsa\_salud.db} & Base de datos SQLite persistida en disco $\rightarrow$ \texttt{data/} \\
\bottomrule
\end{tabularx}
\caption{Trazabilidad de objetos --- DC-09}
\label{tab:col-objetos-dc09}
\end{table}
